\section{Setting up the experiments}
\label{sec:setup}
\textbf{Code:} \code{data.py}, \code{decision.py}, \code{metrics.py}, \code{config.py},
\code{plotting.py}, \code{run\_all.py}.

Before any model could be compared, every experiment had to use the same data, the same splits and
the same way of turning probabilities into actions and costs.

\paragraph{Data and splits.}
The UCI \emph{Spambase} data set has $4{,}601$ emails, $1{,}813$ ($39.4\%$) of them spam, described by
57 word, character and capital-letter statistics. The features are heavily skewed and $77\%$ of all
values are exactly zero. \code{data.py} downloads and checksums the data, makes stratified
$60/20/20$ train/validation/test splits ($2{,}761/921/919$ emails), and fits the standardisation
(optionally after $\log(1+x)$) on the training part only. Every experiment repeats over the same 20
seeds and reports the mean and its standard error. Models are fitted on training data, every
choice (penalty, shrinkage, temperature) is made on validation data, and the test part is used only
to report. \code{data.py} also draws synthetic Gaussian data, whose exact posterior is known, and
resamples a test set within each class to a new spam rate, which Section~\ref{sec:shift} needs.

\paragraph{Decisions and scoring.}
\code{decision.py} turns any loss matrix into the thresholds derived in the next section and applies
them, either to probabilities or directly to log-odds (which avoids rounding a confident score to
exactly $0$ or $1$). A unit test compares it with a brute-force search over the three actions on
$10{,}000$ random posteriors and loss matrices. \code{metrics.py} computes the realized loss, accuracy,
inspection rate and confusion matrices with an inspect column. \code{run\_all.py} regenerates every
figure and table in this report in under a minute, and 101 unit tests cover the shared code and the
models. With data, splits and scoring fixed, the first question was which action to take for an
email once we know how likely it is to be spam.
